\section{Introduction}
\postsec
A foundational goal of machine learning is to develop models that work reliably across a broad range of data distributions. 
ecently, large pre-trained models such as CLIP \cite{radford2021learning} and ALIGN \cite{jia2021scaling} have demonstrated unprecedented robustness to challenging distribution shifts where prior robustness interventions failed to improve performance \cite{taori2020measuring}.
While these results point towards pre-training on large, heterogeneous datasets as a promising direction for increasing robustness, an important caveat is that these robustness improvements occur only in the zero-shot setting, i.e., when the model performs inference without fine-tuning on a specific target distribution \cite{radford2021learning}.


In a concrete application, a zero-shot model can be fine-tuned on extra application-specific data, which often yields large performance gains on the target distribution.
However, Radford \etal \cite{radford2021learning} have shown that current fine-tuning techniques carry a hidden cost: the out-of-distribution accuracy of their fine-tuned CLIP models is often lower than that of the original zero-shot model.
This leads to a natural question: \emph{Can zero-shot models be fine-tuned without reducing out-of-distribution accuracy?}

\begin{figure*}
    \centering
    \includegraphics[width=\textwidth]{new_arxiv_v2/figures/fig1_new.pdf}
    \caption{
    \textbf{(Top left)} Zero-shot CLIP models exhibit high effective robustness and moderate in-distribution accuracy, while standard fine-tuning---either end-to-end or with a linear classifier (final layer)---attains higher ID accuracy and less effective robustness.
    \textbf{(Top right)} Our method linearly interpolates between the zero-shot and fine-tuned models with a mixing coefficient $\alpha \in [0,1]$. \textbf{(Bottom)} On five distribution shifts derived from ImageNet (ImageNetV2, ImageNet-R, ImageNet Sketch, ObjectNet, and ImageNet-A), WiSE-FT improves average OOD accuracy relative to both the zero-shot and fine-tuned models while maintaining or improving ID accuracy.
    }
    \label{fig:fig1}
\end{figure*}

As pre-trained models are becoming a cornerstone of machine learning, techniques for fine-tuning them on downstream applications are increasingly important.
Indeed, the question of robustly fine-tuning pre-trained models has recently also been raised as an open problem by several authors \cite{andreassen2021evolution,bommasani2021opportunities,radford2021learning}. Andreassen \etal \cite{andreassen2021evolution} explored several fine-tuning approaches but found that none yielded models with improved robustness at high accuracy.
Furthermore, Taori \etal \cite{taori2020measuring} demonstrated that no current algorithmic robustness interventions provide consistent gains across the distribution shifts where zero-shot CLIP excels.

In this paper, we conduct an empirical investigation to understand and improve fine-tuning of zero-shot models.
First, we measure how different fine-tuning approaches (last layer vs. end-to-end fine-tuning, hyperparameter changes, etc.) affect the out-of-distribution accuracy of the resulting fine-tuned models.
We find that more aggressive fine-tuning (e.g., using a larger step size) yields larger in-distribution improvements but can also reduce out-of-distribution accuracy by a larger amount.
In addition, out-of-distribution accuracy varies substantially even among models with the same in-distribution accuracy.
Overall, the fine-tuning process is brittle and the best hyperparameters cannot be inferred from in-distribution accuracy alone.

To address these issues, we then introduce a robust way of fine-tuning zero-shot models that achieves the best of both worlds: increased performance out-of-distribution while maintaining or even improving in-distribution accuracy relative to standard fine-tuning. In addition, our method simplifies the choice of hyperparameters in the fine-tuning process.
Our method (Figure \ref{fig:fig1}) has two steps: first, we fine-tune the zero-shot model on application-specific data. Second, we combine the original zero-shot and fine-tuned models by linearly interpolating between their weights, which we refer to as weight-space ensembling. 

Interestingly, linear interpolation in weight space succeeds despite non-linearity in the activation functions of the neural networks.
Weight-space ensembles for fine-tuning (WiSE-FT) substantially improve out-of-distribution accuracy compared to prior work while maintaining high in-distribution performance.
Concretely, on ImageNet \cite{deng2009imagenet} and five of the natural distribution shifts studied by Radford \etal~\cite{radford2021learning}, WiSE-FT applied to standard end-to-end fine-tuning improves out-of-distribution accuracy by 3.9 to 6 percentage points (pp) over prior work while improving the in-distribution accuracy by 1.6 pp.
Relative to the zero-shot model, WiSE-FT improves out-of-distribution accuracy by 1 to 9 pp. Moreover, WiSE-FT improves over a range of alternative approaches such as regularization and evaluating at various points throughout fine-tuning. The robustness comes at no additional computational cost during fine-tuning or inference because the zero-shot and fine-tuned models are ensembled into a single model of the same size.

To understand the robustness gains of WiSE-FT, we empirically analyze ensembling through the lens of distributional robustness. First, we study WiSE-FT when fine-tuning a linear classifier (last layer) as it is amenable to analysis: our procedure is then equivalent to ensembling the outputs of two models, and experiments point towards the complementarity of model predictions as a key property. We illustrate via detailed measurements how the predictions of zero-shot and fine-tuned models are diverse, and that models are more confident on the parts of the test distributions they perform well on.
 \begin{figure*}[t]
    \centering
    \includegraphics[width=\textwidth]{new_arxiv_v2/figures/datasets.pdf}
    \caption{Samples of the class \textit{lemon}, from ImageNet \cite{deng2009imagenet} (in-distribution) and the derived out-of-distribution datasets considered in our main experiments: ImageNet-V2 \cite{pmlr-v97-recht19a}, ImageNet-R \cite{imagenetr}, ImageNet Sketch \cite{imagenetsketch}, ObjectNet \cite{objectnet}, and ImageNet-A \cite{imageneta}.}
    \label{fig:data_samples}
\end{figure*}

For end-to-end fine-tuning, we connect our observations to earlier work on the phenomenology of deep learning. Neyshabur \etal \cite{neyshabur2020being} found that end-to-end fine-tuning the same model twice yielded two different solutions that were connected via a linear path in weight space along which error remains low, known as linear mode connectivity \cite{frankle2020linear}. The authors therefore concluded that the two solutions are in the same basin of the loss landscape. Interestingly, linear interpolation in weight space succeeds despite non-linearity in the activation functions of the neural networks. 
Our observations suggest a similar phenomenon along the path generated by WiSE-FT, but the exact shape of the loss landscape and connection between in- and out-of-distribution error are still an open problem.

In addition to the aforementioned ImageNet distribution shifts, WiSE-FT consistently improves robustness on a diverse set of six further distribution shifts including: (i) geographic shifts in satellite imagery and wildlife recognition (WILDS-FMoW, WILDS-iWildCam) \cite{wilds2021, christie2018functional, beery2021iwildcam}, (ii) reproductions of the popular image classification dataset CIFAR-10 with a distribution shift (CIFAR-10.1 and CIFAR-10.2) \cite{pmlr-v97-recht19a, lu2020harder}, and (iii) datasets with distribution shift induced by temporal perturbations in videos (ImageNet-Vid-Robust and YTBB-Robust) \cite{vidrobust}.



Beyond the robustness perspective, WiSE-FT also improves in-distribution performance compared to standard fine-tuning, reducing the relative error rate by 6 to 62\% on a range of seven datasets: ImageNet, CIFAR-10, CIFAR-100 \cite{krizhevsky2009learning}, Describable Textures \cite{dtd}, Food-101 \cite{food101}, SUN397 \cite{sun397}, and Stanford Cars \cite{cars}.
Even when fine-tuning data is scarce—reflecting many application scenario---we find that WiSE-FT substantially improves in-distribution performance.
When five examples per class are used for fine-tuning, WiSE-FT improves in-distribution performance by 0.3 to 6.1 percentage points compared to the best of the zero-shot and fine-tuned models on the seven aforementioned datasets.
Overall, we find that WiSE-FT is universally applicable in the problems we studied, and we encourage its adoption for fine-tuning zero-shot models.





